\subsection{Límite del pseudodeterminante}\label{sec:D04}
Se usan las definiciones y la convergencia de Hilbert--Schmidt de la Sección~\ref{sec:D03}.
Fuera de la imagen de $E_n$, $\widehat J_n$ es cero.
Por la brecha en dimensión finita,
\[
 I+\widehat J_n\ge aI,\qquad a=\min(1,\gamma)>0.
\]
La convergencia en norma de operadores, consecuencia de la
convergencia Hilbert--Schmidt, da también $I+J\ge aI$.

\paragraph{Logaritmo regularizado.}
Para $B$ autoadjunto Hilbert--Schmidt y $I+B\ge aI>0$, definimos
\[
 \mathcal D(B)=-\int_0^1 t\,
       \operatorname{tr}\bigl[B^2(I+tB)^{-1}\bigr]\,dt
             =\log\det{}_2(I+B).                    \tag{D4.1}
\]
La integral está bien definida: $B^2$ es de clase traza y
$\|(I+tB)^{-1}\|\le\min(1,a)^{-1}$.
En dimensión finita, integrar la identidad escalar
$\log(1+\lambda)-\lambda
=-\int_0^1 t\lambda^2/(1+t\lambda)\,dt$ prueba
\[
 \log\det(I+B)=\operatorname{tr}B+\mathcal D(B).       \tag{D4.2}
\]
La misma identidad espectral define el determinante regularizado
en dimensión infinita, pues sus sumandos son $O(\lambda^2)$ cerca de cero.

\begin{lema}[Continuidad del logaritmo regularizado]\label{lem:continuidad}
Si $B_n\to B$ en Hilbert--Schmidt, todos son autoadjuntos y
$I+B_n\ge aI>0$ uniformemente, entonces
$\mathcal D(B_n)\to\mathcal D(B)$.
\end{lema}
\begin{proof}
La desigualdad de productos Hilbert--Schmidt da
\[
 \|B_n^2-B^2\|_1\le
       (\|B_n\|_2+\|B\|_2)\|B_n-B\|_2\longrightarrow0.
\]
Por la identidad de la resolvente, las inversas $(I+tB_n)^{-1}$
convergen en norma de operadores uniformemente para $0\le t\le1$.
Separar la diferencia de los productos de (D4.1) en dos sumandos,
y usar $|\operatorname{tr}(UV)|\le\|U\|_1\|V\|$,
prueba la convergencia uniforme de los integrandos y la convergencia de sus integrales.
\end{proof}

\begin{teorema}[Límite del pseudodeterminante canónico]\label{thm:pdet}
Con $I_2=\int\nu^2$ y $J=\mathcal K+P_0$ definido en la Sección~\ref{sec:D03},
\[
 \log\operatorname{pdet}S_n
 \longrightarrow C_S:=1+12I_2+\mathcal D(J).           \tag{D4.3}
\]
\end{teorema}
\begin{proof}
El único autovalor cero de $S_n$ es reemplazado por uno en
$I+J_n=S_n+P_n$, por lo que (D4.2) da exactamente
\[
 \log\operatorname{pdet}S_n
       =\operatorname{tr}J_n+\mathcal D(\widehat J_n).
\]
Los ceros adicionales del espectro de $\widehat J_n$ no contribuyen a
$\mathcal D$. La fórmula (D3.2) da el límite de la traza;
(D3.6) y el Lema~\ref{lem:continuidad} dan el segundo límite.
\end{proof}

\paragraph{Observación.}
El término $12I_2$ conserva la traza de una diagonal que desaparece
en norma de Hilbert--Schmidt: no puede obtenerse tomando la traza del operador límite.
La identidad integral (D4.1) sólo usa la positividad en el segmento $I+tJ$;
no presupone $\|J\|<1$.
